\documentclass[11pt]{article}

\usepackage[a4paper,margin=25mm]{geometry}
\usepackage{amsmath,amssymb,mathtools}
\usepackage{bm}
\usepackage{enumitem}
\usepackage{booktabs}
\usepackage{hyperref}

\newcommand{\Cset}{\mathcal C}
\newcommand{\M}{\mathcal M}
\newcommand{\D}{\mathsf D}
\newcommand{\wt}{W}
\newcommand{\Prob}{\Pr}
\newcommand{\supp}{\operatorname{supp}}
\newcommand{\argmin}{\operatorname*{arg\,min}}

\title{Color-Correlated Concatenated MWPM Decoding for the Color Code}
\author{}
\date{}

\begin{document}
\maketitle

\section{Scope and assumptions}

We consider a two-dimensional color code decoded by the concatenated minimum-weight
perfect matching (MWPM) decoder of Lee, Li, and Bartlett.  The Pauli sectors are
treated independently: throughout the derivation, one fixes either the $X$-error
sector or the $Z$-error sector and applies the same construction to the other sector
independently.

Let
\[
    \Cset = \{r,g,b\}
\]
denote the three colors.  A sector-specific detector error model (DEM) is written as
\[
    \M = \{m_i=(q_i,D_i,O_i)\}_{i=1}^{M},
\]
where $q_i$ is the prior probability of the $i$th independent error mechanism,
$D_i$ is the set of detectors flipped by that mechanism, and $O_i$ is the set of
logical observables flipped by it.

For each color $c\in\Cset$, the concatenated decoder constructs two graphlike DEMs:
a $c$-restricted first-stage DEM and a $c$-only second-stage DEM.  We denote these by
\[
    \M_c^{(1)},\qquad \M_c^{(2)}.
\]
The corresponding error probabilities are
\[
    \bm p_c^{(1)},\qquad \bm p_c^{(2)},
\]
and the matching weight of an error mechanism with probability $p$ is
\[
    \ell(p)=\log\frac{1-p}{p}.
\]

\section{Baseline concatenated MWPM}

For a fixed color $c$, let $\{c_1,c_2\}=\Cset\setminus\{c\}$.  The first MWPM is
performed on the $c$-restricted DEM using the violated detectors of colors $c_1$ and
$c_2$.  Its solution is denoted by
\[
    x_c^{(1)}
    =
    \operatorname{MWPM}
    \left(
        \sigma^{(c_1)}\cup\sigma^{(c_2)};
        \M_c^{(1)}
    \right).
\]
The selected first-stage mechanisms are converted into virtual detectors and combined
with the physical $c$-colored detector outcomes.  The second MWPM returns
\[
    x_c^{(2)}
    =
    \operatorname{MWPM}
    \left(
        \sigma^{(c)}\cup V(x_c^{(1)});
        \M_c^{(2)}
    \right).
\]
The resulting second-stage correction is mapped back to the original DEM ordering:
\[
    e_c^{(0)} = R_c x_c^{(2)},
\]
where $R_c$ denotes the existing error-map operation implemented by the decoder.

The baseline solution weight, using the ordinary unmodified second-stage prior of
color $c$, is
\[
    W_c^{(0)}
    =
    \sum_j x_{c,j}^{(2)}
    \log\frac{1-p_{c,j}^{(2)}}{p_{c,j}^{(2)}}.
\]
The ordinary concatenated decoder obtains the three candidates
\[
    e_r^{(0)},\qquad e_g^{(0)},\qquad e_b^{(0)}
\]
and chooses the least-weight one.

\section{Source provenance of decomposed mechanisms}

Each error mechanism $a$ in a decomposed DEM is derived from one or more mechanisms
of the original DEM.  Let
\[
    \pi(a)\subseteq\{1,\ldots,M\}
\]
denote its source set.  In the current implementation, this information is represented
by the decomposition error-map matrices.

For the standard decomposition in which the source mechanisms appear with unit
multiplicity, the marginal probability of a decomposed mechanism $a$ is
\[
    p_a
    =
    \frac{
        1-\prod_{i\in\pi(a)}(1-2q_i)
    }{2}.
\]
This is the probability that an odd number of the independent source mechanisms in
$\pi(a)$ occurs.

\section{Color-correlated prior update}

The purpose of the color-correlated step is to use a candidate correction obtained
from one color as evidence for reweighting the two-stage matching problem of a
different target color.

Let $G\subseteq\{1,\ldots,M\}$ be the set of original DEM mechanisms selected by a
guide correction.  For a target decomposed mechanism $a$, consider conditioning on a
guide source mechanism $i\in G$ having occurred.

If $i\notin\pi(a)$, the target mechanism is independent of that guide source and
\[
    p_{a\mid i}=p_a.
\]
If $i\in\pi(a)$, conditioning on $i$ being active flips the parity condition for the
remaining sources, yielding
\[
    p_{a\mid i}
    =
    \frac{
        1+\displaystyle\prod_{j\in\pi(a)\setminus\{i\}}(1-2q_j)
    }{2}.
\]
The reweighted prior is defined by the same ``take the strongest implied
probability'' rule used in correlated matching:
\[
    \widetilde p_a^{(G)}
    =
    \max\left(
        p_a,\,
        \max_{i\in G}p_{a\mid i}
    \right).
\]
Numerically, if the conditional probability is exactly one, it should be clipped to
$1-\varepsilon$ before converting it to a finite log-likelihood weight:
\[
    \widetilde\ell_a^{(G)}
    =
    \log\frac{1-\widetilde p_a^{(G)}}{\widetilde p_a^{(G)}}.
\]

The update is applied to \emph{both} stages of the target-color concatenated decoder:
\[
    \M_c^{(1)}
    \longrightarrow
    \widetilde\M_{c,G}^{(1)},
    \qquad
    \M_c^{(2)}
    \longrightarrow
    \widetilde\M_{c,G}^{(2)}.
\]
The first-stage MWPM must therefore be rerun; reweighting only the second stage is not
equivalent to rerunning the concatenated decoder because the first-stage solution
determines the virtual syndrome used by the second stage.

\section{Guide sets}

For the baseline color-$d$ solution, define
\[
    G_d
    =
    \supp(e_d^{(0)})
    =
    \{i:e_{d,i}^{(0)}=1\}.
\]
For the combined guide from two colors $d_1$ and $d_2$, the operation is a Boolean
OR, i.e.\ a set union rather than an XOR:
\[
    G_{d_1\lor d_2}
    =
    G_{d_1}\cup G_{d_2}.
\]

For a target color $c$ with other colors $\{d_1,d_2\}=\Cset\setminus\{c\}$, run the
full two-stage concatenated decoder three additional times:
\begin{align}
    e_c^{[d_1]}
    &=
    \D_c\!\left(
        \sigma;
        \widetilde\M_{c,G_{d_1}}^{(1)},
        \widetilde\M_{c,G_{d_1}}^{(2)}
    \right),\\
    e_c^{[d_2]}
    &=
    \D_c\!\left(
        \sigma;
        \widetilde\M_{c,G_{d_2}}^{(1)},
        \widetilde\M_{c,G_{d_2}}^{(2)}
    \right),\\
    e_c^{[d_1\lor d_2]}
    &=
    \D_c\!\left(
        \sigma;
        \widetilde\M_{c,G_{d_1\lor d_2}}^{(1)},
        \widetilde\M_{c,G_{d_1\lor d_2}}^{(2)}
    \right).
\end{align}

\section{Twelve-candidate set}

The full candidate set is
\begin{align}
\mathcal E
={}&
\{
e_r^{(0)},e_g^{(0)},e_b^{(0)}
\}
\nonumber\\
&\cup
\{
e_r^{[g]},e_r^{[b]},e_r^{[g\lor b]}
\}
\nonumber\\
&\cup
\{
e_g^{[r]},e_g^{[b]},e_g^{[r\lor b]}
\}
\nonumber\\
&\cup
\{
e_b^{[r]},e_b^{[g]},e_b^{[r\lor g]}
\}.
\end{align}
Thus, before deduplication,
\[
    |\mathcal E|=12.
\]
The first version of the algorithm is non-iterative: only the three baseline solutions
are used as guides.  Correlated candidates are not recursively used to generate
further candidates.

\section{Final scoring under a common prior}

The weights produced during the nine correlated reruns are evaluated under different
reweighted priors and therefore cannot be compared directly.  Reweighting is used
only for \emph{candidate generation}.

For each candidate generated for target color $c$, retain its native second-stage
binary solution $x_c$.  Recompute its selection weight with the ordinary,
unmodified second-stage probabilities of the same color:
\[
    W_0(x_c)
    =
    \sum_j x_{c,j}
    \log\frac{1-p_{c,j}^{(2)}}{p_{c,j}^{(2)}}.
\]
This definition deliberately matches the existing concatenated decoder's weight
semantics.  The final output is
\[
    e^\star
    =
    \argmin_{e\in\mathcal E} W_0(e).
\]

Hence the algorithm has the structure
\[
    \text{baseline decoding}
    \;\longrightarrow\;
    \text{cross-color conditional reweighting}
    \;\longrightarrow\;
    \text{nine additional candidates}
    \;\longrightarrow\;
    \text{common-prior rescoring}
    \;\longrightarrow\;
    \text{minimum-weight output}.
\]

\section{Comparative decoding}

If comparative decoding is enabled, the entire 12-candidate construction is performed
independently inside each fixed logical class $\lambda$.  Let
\[
    W_{\lambda,k},
    \qquad k=1,\ldots,12,
\]
denote the common-prior weights of the 12 candidates in logical class $\lambda$.
The class weight is
\[
    W_\lambda=\min_k W_{\lambda,k}.
\]
The final logical class is the class with minimum $W_\lambda$, and the existing
logical-gap definition is preserved:
\[
    \Delta
    =
    W_{(2)}-W_{(1)},
\]
where $W_{(1)}\le W_{(2)}\le\cdots$ are the sorted class minima.

\section{Implementation invariants}

The implementation should preserve the following invariants.

\begin{enumerate}[leftmargin=2em]
    \item When color-correlated decoding is disabled, the decoder must be bit-for-bit
    compatible with the current concatenated decoder for the same inputs and random
    seed.
    \item Every returned final correction remains represented in the original DEM
    ordering, as in the current implementation.
    \item The public \texttt{weights} output remains a weight under the ordinary
    unmodified prior, never under a correlated reweighted prior.
    \item Existing logical-gap computation is performed from the minimum
    common-prior candidate weight in each logical class.
    \item The reported \texttt{best\_colors} value continues to denote the target
    color $r/g/b$, not the index of one of the 12 candidate-generation branches.
    \item Boolean OR of two guide corrections means union of selected error mechanisms,
    not symmetric difference.
    \item Existing correction-based metrics can operate on the final mapped correction
    without modification.
    \item Any soft-output quantity that depends on the internal matching growth state
    must not be silently reused for a candidate produced under a different reweighted
    graph.  Such combinations should either be recomputed consistently or rejected
    explicitly.
\end{enumerate}

\begin{thebibliography}{9}

\bibitem{LeeLiBartlett}
S.-H. Lee, A. Li, and S. D. Bartlett,
``Color code decoder with improved scaling for correcting circuit-level noise,''
\emph{Quantum} (2025).

\bibitem{ShuttyNewmanVillalonga}
N. Shutty, M. Newman, and B. Villalonga,
``Efficient near-optimal decoding of the surface code through ensembling,''
arXiv:2401.12434.

\end{thebibliography}

\end{document}
